% Sol, step-15 manuscript contribution. New prose; locked outline unchanged.
% Requires amsmath and booktabs. Citation keys belong to phase-3 references.bib.
% Sections may be integrated separately. Numbers are finite-cohort descriptions.

\section{Study design and economic measurement}
\label{sec:sol-economic-methods}

\subsection{The cohort and the object of evaluation}

We study twenty completed four-player games, comprising eighty player-games, 50,918 events, 6,989 decisions and 8,679 response attempts. Five seeds each have four cyclic seat rotations. Sol, Gemini, Opus and Grok each occupy seats 1--4 five times. The evaluated participant is the configured agent, including its model, inference settings, observations, memory and response-resolution policy. This complete-system interpretation follows existing arguments for agent evaluation and is especially consequential when a retry changes the response condition~\cite{kapoor2026hal,zhang2026dece}. An attempt is one recorded response within a decision. We write G$n$ d$k$ for zero-based decision $k$ in game $n$ (Appendix~\ref{app:cohort}). The engine applies one resolved action for that decision; several decisions can occur during one player's turn.

The finite cohort is our population of description. Seat balance does not make the games independent draws from a common distribution of opportunities. Decisions within a game share an evolving board, and rotations within a seed block share initial random sequences. Purchases, jail choices and eliminations alter how those sequences are consumed. Two saved treatment fingerprints also divide games 1--4 from games 5--20. The inspected source diff changes campaign-concurrency plumbing and other application files without changing engine, prompt, resolver or model-client files. This inspection does not identify an effect, or establish the absence of one. We therefore report observed outcomes and paths without treating thousands of decisions as independent evidence of an intrinsic model ranking.

Our analysis combines whole trajectories with narrower opportunity-conditioned summaries. This follows the general motivation for process-sensitive agent evaluation~\cite{ma2024agentboard}, but requires explicit economic units. A deed acquisition changes ownership of one property. A group completion occurs when a player acquires the last deed needed for a color group. A trade is an atomic settlement between two players, with its transfers and charges evaluated together. We distinguish the original thread initiator, the author of the final offer and its accepting player. These roles need not coincide. We record completion at stable transaction boundaries so that an intermediate transfer inside a swap cannot create a fictitious opportunity.

\subsection{Opportunities, expenditures and observed receipts}

For each trade-completed group, the first management window begins at the owner's first actual optional-management decision after settlement and includes further management decisions in that same global turn. We distinguish construction legal on entry, construction becoming legal later in that window, and construction actually applied. Jail and other intervening phases can postpone management. Financing choices within the window can create a build opportunity. Thus, conditioning on entry buildability answers a different question from following every acquisition. It does not remove the consequences of earlier play.

We measure negotiation initiation using two denominators. For configured agent $m$, let $\mathcal{T}_m$ contain its own turns with at least one menu permitting a proposal, and let $\mathcal{P}_m\subseteq\mathcal{T}_m$ contain turns with an applied proposal. If $N_m$ is its number of proposals, then
\begin{equation}
 q_m=\frac{|\mathcal{P}_m|}{|\mathcal{T}_m|},
 \qquad
 i_m=\frac{N_m}{|\mathcal{P}_m|}.
 \label{eq:sol-proposals}
\end{equation}
The first quantity describes proposing-turn propensity; the second describes proposals per proposing turn. Neither is bargaining quality. A denominator of eligible decisions would also count repeated management requests caused by earlier mortgages or purchases, rather than additional turns. We retain that count for auditability without equating it with independent access.

A spending unit $u=(g,p,c)$ joins game $g$, player $p$ and color group $c$ with positive own expenditure on houses or hotels. Including hotels prevents the omission of an owner who inherits four-house deeds and purchases only hotel conversions. Let $\tau_u$ be the first own construction event and $T_g$ the game's terminal event. We report
\begin{equation}
 K_u=\sum_{e\in\mathcal{B}_u}\operatorname{cost}(e),
 \qquad
 R_u^{\mathrm{post}}=
 \sum_{\substack{e\in\mathcal{R}_u\\\tau_u\leq\operatorname{seq}(e)\leq T_g}}
 \operatorname{paid}(e).
 \label{eq:sol-receipts}
\end{equation}
Here $\mathcal{B}_u$ contains that owner's construction events and $\mathcal{R}_u$ contains paid rent on that owner's developed deeds in the group. We also report whole-game developed receipts, without the $\tau_u$ filter. All amounts are game dollars. Acquisition costs, redemption charges, building salvage, mortgage proceeds and estate transfers remain separate flows. An unpaid terminal rent demand is not a paid receipt.

These quantities do not estimate incremental investment returns. The owner can inherit improvements, and subsequent rent depends on opponents' movement and survival. Removing receipts before first own spending fixes a timing mismatch while leaving inherited productive assets in the later window. Observation ends at ownership loss or game end. A late group and an early group consequently have different exposure. We preserve the sequence of acquisitions and expenditures instead of treating $R_u^{\mathrm{post}}/K_u$ as a comparable return on investment. Similar Monopoly investment questions appear in FinBoardBench; our event accounting addresses realized flows and legal financing rather than a universal valuation of the acquired assets~\cite{hu2026finboardbench}.

Ownership, physical development, cash and bank stock are reconstructed from events. The retained audit reconciles these quantities with all 6,989 initial decision payloads and checks construction and salvage against cash-change events. We distinguish a group's completion route from the acquisition routes of its constituent deeds. A group can contain inherited deeds but become complete through a later purchase or trade. No spending unit in this cohort has a second complete-ownership interval at the observed decision or terminal boundaries.

\subsection{Local legal-sequence tests}

To test a financing or construction claim, we restore the actual historical decision prefix, check its event agreement, and execute a specified sequence of the acting player's own legal actions through the engine. We stop at a declared local endpoint, such as completing a proposed construction multiset or paying the current debt. A test cannot borrow another player's management opportunity. Conditional acceptance, altered stock or other diagnostic changes must be stated separately. Legal-action interfaces and replay already occur in related economic benchmarks~\cite{wu2026agentmonopoly,zhang2026dece}; the present use is to adjudicate concrete claims about recorded trajectories.

For a debtor with cash $C$, a generous nominal upper bound is
\begin{equation}
 U=C+V_{\mathrm{all\ building\ salvage}}
       +V_{\mathrm{all\ clear\ deed\ mortgages}}.
 \label{eq:sol-upper}
\end{equation}
The mortgage term counts each unmortgaged deed once, including railroads and utilities and collateral that clearing buildings could unlock; already-mortgaged deeds contribute nothing. If $U$ is below debt $D$, own-asset liquidation cannot pay $D$. If $U\geq D$, a legal sequence must still establish access. A hotel downgrade needs bank houses, and the even-selling rule can prevent reaching nominal salvage. Conversely, a current menu can understate finance because mortgages become available only after sales. A tested payment route establishes current-debt feasibility, not future survival or the value of surrender. Once a branch changes the state, the historical policy suffix does not identify how any player would continue.


\subsection{Claims, interventions and certificates}
\label{sec:claim-protocol}

The audit unit is a claim $Q=(h,p,\mathcal{I},E)$: a hashed historical prefix $h$, the acting player $p$, an admissible intervention class $\mathcal{I}$, and an endpoint predicate $E$. This specification prevents ``could pay'' from silently meaning ``could obtain help'', or ``could build'' from borrowing an earlier management phase. A local alternative changes only the permitted actions from the restored prefix. It does not retain a historical future after that change.

\begin{table}[t]
\centering\small
\begin{tabularx}{\linewidth}{@{}L{0.20\linewidth}Y@{}}
\toprule
Audit step & Required evidence and stopping rule \\
\midrule
Fix the claim & Record actor, decision, source hashes, allowed actions, endpoint and assumptions before executing the diagnostic. Disclose retrospective case selection.\\
Restore the prefix & Match the pending decision and actor, and every preceding canonical event. Restore through original actions, including random and card state, rather than edit a displayed snapshot.\\
Establish feasibility & Execute a permitted sequence through the unchanged engine. Preserve actions and events; stop when the endpoint is met. One witness proves existence, not optimality.\\
Establish impossibility & Use an admissible bound covering all allowed inflows, or exhaust the permitted state graph. A failed heuristic or timed-out search returns unresolved.\\
Limit the inference & Separate reachability from submitted encoding, future survival, predicted agreement and model choice. Reconstruct delivery and recipient response separately for social claims.\\
\bottomrule
\end{tabularx}
\caption{Historical-claim audit procedure. The endpoint and intervention class determine the certificate's meaning. Engine execution tests feasibility, not belief, attention or persuasion.}
\label{tab:claim-protocol}
\end{table}

We implement this procedure for two complete, retrospectively selected populations: every elimination and every engine-defaulted construction decision. Selection comes from canonical bankruptcy events and resolution records before the new runner executes its diagnostics. For liquidation, $\mathcal{I}$ permits own legal sales and mortgages against the existing debt, without help, trades, construction or future random moves. Equation~\ref{eq:sol-upper} is admissible because each sale consumes finite salvage and each mortgage consumes one clear-deed capacity. Ignoring stock and even-selling restrictions enlarges that budget. Where the bound does not settle the claim, we retain a payment witness or a complete financing graph.

For construction, $\mathcal{I}$ permits splitting count-valued operations into atomic purchases and changing their order, preserving every deed/kind count. It supplies no additional cash, stock or action slot. We test the submitted encoding separately. A reachable target is therefore compatible with an invalid attempted action. Our constructive ordering is not a complete solver for arbitrary targets; failure with adequate resources would remain unresolved. No such residual occurs here: all well-formed unsuccessful targets exceed current cash.

The portable artifact contains the fixed 126-case manifest, original action streams, per-event reference digests, settings, an unchanged engine and board copy, pinned dependencies, witnesses and negative certificates. It reproduces all 50,918 original events and all 126 audited prefixes. Nine adversarial tests reject altered inputs and invalid certificates. One new test preserves a reviewer-found counterexample: the engine accepts a bundled building sale that an atomic search misses. The checker therefore refuses a graph-based impossibility certificate whenever a visited state offers building sales. The sole historical graph certificate has no sale action and is exhaustive over its mortgages. This correction illustrates why engine acceptance of positive witnesses is easier to establish than completeness of negative searches. The artifact does not measure annotation accuracy against a human standard or establish transfer to another environment.


\subsection{Reconstructing conduct and recipient response}
\label{sec:conduct-method}

Social analysis begins with a proposition or commitment in its chronological context, rather than a global label attached to a colorful exchange. We link the speaker's observation, public message, selected action, resulting events and any saved private note or provider explanation. We then identify which recipient prompts contain the relevant statement, both earlier in the interaction and at the decision where reliance is proposed. A recipient may have acted on earlier exposure even if the message has since left its current window. Conversely, a message's presence in a canonical event does not prove delivery to that recipient.

We adjudicate five questions separately. Is the stated proposition false or misleading under a specified reading and the speaker's available information? Does another recorded channel contradict it? Was the proposition delivered? Does the recipient's subsequent record endorse that proposition or act in documented reliance on it? Does the applicable instruction or agreement prohibit the conduct? These dimensions are not an escalating deception score. An observed breach can remain a breach after a promise scrolls out of memory; missing context instead limits explanations based on knowing violation or forgetting. A trade can settle while the recipient expressly rejects its public premise. A statement can be false without evidence that the speaker knew it was false.

Provider explanations and private notes are recorded text, not direct measurements of internal computation. We preserve agreement and disagreement among channels instead of assigning one privileged access to intent. This is especially necessary when a received tool call conflicts with the accompanying content or when an explanation corrects itself before the selected action. Work on explanation faithfulness motivates this caution, but cannot determine the mechanism in any particular record~\cite{turpin2023unfaithful,lanham2023faithfulness}. Mechanically verifiable propositions, such as ownership, charges or the feasibility of a stated action sequence, can still be checked independently.

Detection opportunities differ by configuration. Private notes are largely absent under Opus's action-only condition, provider explanations vary in availability, and content/tool conflicts are especially observable in Grok's records. We identify public actions and statements separately from channel-qualified case evidence. A rate among decisions with readable explanations would remain a selected-subset rate: refusal, recovery and game progression affect that selection. No adjustment by channel count alone establishes comparable latent deception or intent prevalence.

\subsection{Review, counterexamples and claim selection}
\label{sec:review}

The corpus received a complete public-scope pass and a complete pass over all available saved channels and attempts, with unavailable or encrypted material marked as such. The extraction stage produced 441 public-scope candidates and 152 final full-channel candidates, alongside twenty-one historical checkpoints. Four coauthor review streams each considered all 614 catalog entries, all twenty trajectories and a common sample of 240 original decision bundles. The catalog contains overlapping observations, not 614 independent episodes. The common sample contains three decisions per player per game from a public-unflagged pool; it serves as an audit of ordinary and missed cases, not a prevalence estimator.

Initial judgments, subsequent source checks and peer-informed corrections are retained separately. This is AI-assisted extraction and adversarial review, not independent human annotation or a validated automated judge. Reading another reviewer's correct result is not counted as an independent derivation. Mechanical disagreements are resolved against source events, accepted transaction payloads and engine-executed prefixes where applicable. Semantic disagreements retain the competing readings if the record does not settle them.

The mechanical populations are censuses; the semantic cases are selected contrasts. We use the latter as counterexamples to specified sufficiency rules, such as ``settlement implies endorsement of the false premise'', or to delimit an unresolved interpretation. One checked counterexample can refute such a rule but cannot establish prevalence. We also preserve analyst errors. Interpreting an \texttt{unmortgageable} menu field as forbidden actions fabricated a financing restriction; treating a sale menu as total capacity omitted collateral unlocked later. These are evaluator errors, not additional agent failures.

